\documentclass[aspectratio=169]{beamer}

% =====================================================================
% ART DECK 07 -- THE BLOOM. Starts in monochrome silver and gains
% colour gradually, slide by slide, until the finale is full neon.
% One saturation knob per slide drives every accent: ncyan!sat!silver
% (0 = silver, 100 = full hue). Two pdflatex passes (overlay).
% =====================================================================

\usepackage{tikz}
\usepackage[T1]{fontenc}
\usetikzlibrary{shadings, fadings, shapes.geometric, shapes.misc,
  decorations.pathmorphing, calc, backgrounds, positioning, arrows.meta,
  shapes.symbols}

\usetheme{default}
\setbeamertemplate{navigation symbols}{}
\setbeamertemplate{footline}{}

% ---- monochrome reference + full vibrant hues -----------------------------
\definecolor{void}{HTML}{08090C}
\definecolor{nightgrey}{HTML}{16181E}
\definecolor{silver}{HTML}{CCD4DE}
\definecolor{paper}{HTML}{EAF2FB}
\definecolor{ncyan}{HTML}{38E8FF}
\definecolor{nmag}{HTML}{FF49B0}
\definecolor{npink}{HTML}{FF77D6}
\definecolor{nlime}{HTML}{B9FF45}
\definecolor{namber}{HTML}{FFC24A}
\definecolor{gold}{HTML}{FFD56B}
\definecolor{eblue}{HTML}{4F7BFF}
\definecolor{nviolet}{HTML}{A05CFF}
\definecolor{nteal}{HTML}{19E6CE}
\definecolor{ngreen}{HTML}{3CE89A}
\definecolor{nred}{HTML}{FF5468}

\tikzfading[name=rglow,  inner color=transparent!0,  outer color=transparent!100]
\tikzfading[name=fadeup, top color=transparent!0,    bottom color=transparent!100]

\pgfmathsetseed{707}

% saturation knob: 0 = silver/monochrome, 100 = full colour
\newcommand{\setsat}[1]{%
  \colorlet{wcyan}{ncyan!#1!silver}%
  \colorlet{wmag}{nmag!#1!silver}%
  \colorlet{wpink}{npink!#1!silver}%
  \colorlet{wlime}{nlime!#1!silver}%
  \colorlet{wamber}{namber!#1!silver}%
  \colorlet{wgold}{gold!#1!silver}%
  \colorlet{weblue}{eblue!#1!silver}%
  \colorlet{wviolet}{nviolet!#1!silver}%
  \colorlet{wteal}{nteal!#1!silver}%
  \colorlet{wgreen}{ngreen!#1!silver}%
  \colorlet{wred}{nred!#1!silver}%
  \colorlet{wroyal}{eblue!#1!silver}}

\newcommand{\art}[1]{%
  \begin{frame}[plain]
  \begin{tikzpicture}[remember picture, overlay,
      shift={(current page.south west)}, x=1cm, y=1cm]
    \fill[void] (0,0) rectangle (16,9);
    #1
  \end{tikzpicture}
  \end{frame}}

\newcommand{\halo}[4]{\fill[#3, opacity=#4, path fading=rglow] #1 circle (#2);}

\newcommand{\spark}[3]{%
  \fill[#3, opacity=0.55, path fading=rglow] #1 circle ({#2*7});
  \fill[#3, opacity=0.85, path fading=rglow] #1 circle ({#2*2.8});
  \begin{scope}[blend mode=screen]
    \fill[white, opacity=0.55, path fading=rglow] #1 ellipse ({#2*9} and {#2*0.45});
    \fill[white, opacity=0.55, path fading=rglow] #1 ellipse ({#2*0.45} and {#2*9});
  \end{scope}
  \fill[white] #1 circle (#2);}

\newcommand{\neon}[3]{%
  \draw[#3, line width=7pt, opacity=0.10] #1 -- #2;
  \draw[#3, line width=3pt, opacity=0.30] #1 -- #2;
  \draw[white, line width=0.7pt, opacity=0.80] #1 -- #2;}

\newcommand{\dust}[1]{%
  \foreach \dk in {1,...,#1}{%
     \pgfmathsetmacro{\sx}{rnd*16}\pgfmathsetmacro{\sy}{rnd*9}%
     \pgfmathsetmacro{\sr}{0.006+rnd*0.020}\pgfmathsetmacro{\so}{0.10+rnd*0.45}%
     \fill[white, opacity=\so] (\sx,\sy) circle (\sr);}}

\newcommand{\bokeh}[1]{%
  \begin{scope}[blend mode=screen]
  \foreach \bk in {1,...,#1}{%
     \pgfmathsetmacro{\bx}{rnd*16}\pgfmathsetmacro{\by}{rnd*9}%
     \pgfmathsetmacro{\br}{0.05+rnd*0.50}\pgfmathsetmacro{\bo}{0.08+rnd*0.32}%
     \pgfmathtruncatemacro{\ci}{rnd*4.999}%
     \ifcase\ci \def\bc{wcyan}\or\def\bc{wviolet}\or\def\bc{wmag}\or
        \def\bc{weblue}\or\def\bc{wteal}\fi
     \fill[\bc, opacity=\bo, path fading=rglow] (\bx,\by) circle (\br);}
  \end{scope}}

\newcommand{\rays}[3]{%
  \begin{scope}[blend mode=screen]
  \foreach \rk in {1,...,#3}{%
     \pgfmathsetmacro{\ang}{rnd*360}\pgfmathsetmacro{\len}{4.5+rnd*7}%
     \pgfmathsetmacro{\wd}{0.4+rnd*1.4}%
     \fill[#2, opacity=0.05, path fading=rglow]
        #1 -- ($#1+(\ang-\wd:\len)$) -- ($#1+(\ang+\wd:\len)$) -- cycle;}
  \end{scope}}

\newcommand{\glowtext}[3]{%
  \begin{scope}[blend mode=screen]
     \fill[#2, opacity=0.50, path fading=rglow] #1 ellipse (5.2 and 1.5);
  \end{scope}
  \node at #1 {#3};}

% nebula using the working (ramped) colours over a neutral ambient
\newcommand{\nebula}{%
  \begin{scope}[blend mode=screen]
     \fill[nightgrey,opacity=0.95, path fading=rglow] (8,4.5) circle (9.5);
     \fill[wviolet, opacity=0.34, path fading=rglow] (5.5,6) circle (6);
     \fill[wmag,    opacity=0.24, path fading=rglow] (11,6) circle (4.8);
     \fill[wcyan,   opacity=0.22, path fading=rglow] (9,3) circle (5.5);
  \end{scope}}

% small saturation meter, bottom-right, so the bloom is legible
\newcommand{\meter}[1]{%
  \draw[silver, opacity=0.35, line width=0.6pt] (13.7,0.45) rectangle (15.5,0.62);
  \fill[wcyan, opacity=0.8] (13.7,0.45) rectangle ({13.7+0.018*#1},0.62);
  \node[anchor=east, text=silver, opacity=0.4, font=\tiny] at (13.6,0.535) {colour};}

\begin{document}

% =====================================================================
% 1. TITLE  -- monochrome
% =====================================================================
\art{
  \setsat{0}
  \nebula \dust{220}\bokeh{55}
  \coordinate (n1) at (3.6,7.0); \coordinate (n2) at (6.6,7.9);
  \coordinate (n3) at (10.2,7.1);\coordinate (n4) at (2.8,4.9);
  \coordinate (n5) at (6.2,5.3); \coordinate (n6) at (10.9,5.1);
  \coordinate (n7) at (13.0,6.3);\coordinate (n8) at (8.3,3.7);
  \coordinate (n9) at (4.9,3.2);
  \foreach \a/\b in {n1/n2,n2/n3,n1/n4,n4/n5,n5/n6,n3/n6,n2/n5,n6/n7,
     n5/n8,n8/n9,n4/n9,n9/n5}{\neon{(\a)}{(\b)}{silver}}
  \foreach \n/\s in {n1/0.075,n2/0.06,n3/0.07,n4/0.055,n5/0.085,
     n6/0.06,n7/0.05,n8/0.06,n9/0.07}{\spark{(\n)}{\s}{silver}}
  \begin{scope}[blend mode=screen]
     \fill[silver, opacity=0.35, path fading=rglow] (5.7,1.85) ellipse (6 and 1.5);
  \end{scope}
  \node[anchor=west, text=white] at (1.3,2.15)
    {\fontsize{46}{50}\selectfont\textbf{Rotor in Rust}};
  \node[anchor=west, text=silver, opacity=0.9] at (1.4,1.2)
    {\large the same story --- told until it finds its colour};
  \node[anchor=west, text=silver, opacity=0.55] at (1.4,0.5)
    {\footnotesize Jasmin Begi\'{c}\; $\cdot$\; Daniel Wassie};
}

% =====================================================================
% 2. THE PROBLEM
% =====================================================================
\art{
  \setsat{7}
  \begin{scope}[blend mode=screen]
     \fill[nightgrey, opacity=0.9, path fading=rglow] (8,4.5) circle (9);
     \fill[wmag, opacity=0.5, path fading=rglow] (11.3,5.7) circle (4.5);\end{scope}
  \foreach \x in {0,...,37}{\foreach \y in {0,...,16}{
     \pgfmathsetmacro{\jx}{\x*0.40+0.6+rnd*0.05}
     \pgfmathsetmacro{\jy}{\y*0.40+1.5+rnd*0.05}
     \fill[weblue, opacity={0.28+rnd*0.30}] (\jx,\jy) circle (0.045);}}
  \dust{110}
  \rays{(11.3,5.7)}{wpink}{24}\spark{(11.3,5.7)}{0.12}{wmag}
  \node[text=white] at (8,8.4) {\Large the one input that crashes \textit{hides} among};
  \glowtext{(8,0.95)}{wcyan}{\fontsize{34}{36}\selectfont\textbf{$2^{64}$ \normalsize\textcolor{silver}{possibilities}}}
  \meter{7}
}

% =====================================================================
% 3. VERIFICATION
% =====================================================================
\art{
  \setsat{15}
  \begin{scope}[blend mode=screen]
     \fill[nightgrey, opacity=0.9, path fading=rglow] (8,4.5) circle (9);
     \fill[weblue, opacity=0.4, path fading=rglow] (8,7) circle (7);\end{scope}
  \foreach \x in {0,...,37}{\foreach \y in {0,...,12}{
     \pgfmathsetmacro{\jx}{\x*0.40+0.6}\pgfmathsetmacro{\jy}{\y*0.40+1.6}
     \fill[wcyan, opacity={0.20+rnd*0.5}] (\jx,\jy) circle (0.05);}}
  \begin{scope}[blend mode=screen]
     \fill[wcyan, opacity=0.20, path fading=rglow] (8,8.4) -- (0.8,1.3) -- (15.2,1.3) -- cycle;
     \fill[white, opacity=0.28, path fading=rglow] (8,8.4) -- (5.4,1.3) -- (10.6,1.3) -- cycle;
  \end{scope}
  \spark{(8,8.3)}{0.12}{wcyan}
  \node[text=white] at (8,0.75) {\Large one question covers \textbf{every} input at once};
  \meter{15}
}

% =====================================================================
% 4. THE PIPELINE
% =====================================================================
\art{
  \setsat{26}
  \begin{scope}[blend mode=screen]
     \fill[nightgrey, opacity=0.9, path fading=rglow] (8,4.6) circle (9);\end{scope}
  \dust{60}
  \def\sy{4.6}
  \draw[wcyan, line width=10pt, opacity=0.10] (2.2,\sy) -- (13.4,\sy);
  \draw[wviolet,line width=4pt,  opacity=0.22] (2.2,\sy) -- (13.4,\sy);
  \foreach \x/\lbl/\c in {2.2/binary/weblue, 6.0/model/wcyan,
        9.8/solver/wviolet, 13.4/answer/wlime}{
     \fill[\c, opacity=0.50, path fading=rglow] (\x,\sy) circle (1.6);
     \shade[inner color=white, outer color=\c!55!void] (\x,\sy) circle (0.7);
     \draw[\c, line width=1.4pt] (\x,\sy) circle (0.7);
     \node[text=void, font=\footnotesize\bfseries] at (\x,\sy) {\lbl};}
  \node[text=white] at (8,8.0) {\Large from a compiled binary to a readable answer};
  \node[text=wcyan, font=\footnotesize] at (6.0,3.4) {rotor};
  \node[text=wlime, font=\footnotesize] at (13.4,3.4){visualiser};
  \meter{26}
}

% =====================================================================
% 5. SPEED
% =====================================================================
\art{
  \setsat{37}
  \begin{scope}[blend mode=screen]
     \fill[nightgrey, opacity=0.9, path fading=rglow] (8,4.5) circle (10);
     \fill[wamber, opacity=0.25, path fading=rglow] (13.4,4.9) circle (4.5);\end{scope}
  \dust{110}
  \begin{scope}[blend mode=screen]
  \foreach \mk in {1,...,38}{
     \pgfmathsetmacro{\my}{1.2+rnd*7.2}\pgfmathsetmacro{\mx}{1+rnd*9}
     \pgfmathsetmacro{\ml}{0.4+rnd*2.2}\pgfmathsetmacro{\mo}{0.05+rnd*0.16}
     \fill[wcyan, opacity=\mo, path fading=rglow] (\mx,\my) ellipse ({\ml} and 0.02);}
  \end{scope}
  \begin{scope}[blend mode=screen]
     \fill[wamber, opacity=0.55, path fading=rglow]
        (1.5,2.7) .. controls (6,3.4) and (10,4.2) .. (13.4,4.9)
        .. controls (10,3.7) and (6,3.0) .. (1.5,2.5) -- cycle;
     \fill[wgold, opacity=0.5, path fading=rglow]
        (4.5,3.3) .. controls (8,3.9) and (11,4.4) .. (13.4,4.9)
        .. controls (11,4.1) and (8,3.6) .. (4.5,3.1) -- cycle;
  \end{scope}
  \rays{(13.5,4.9)}{wgold}{22}\spark{(13.5,4.9)}{0.16}{wamber}
  \node[text=white] at (5.2,7.2) {\fontsize{70}{72}\selectfont\textbf{1000$\times$}};
  \node[text=silver, opacity=0.9] at (5.4,5.7) {\large faster model generation};
  \node[text=silver, opacity=0.75] at (4.2,1.3)
    {\footnotesize 139\,s $\rightarrow$ 0.1\,s \qquad 428\,MB $\rightarrow$ 20\,MB};
  \meter{37}
}

% =====================================================================
% 6. THE COUNT
% =====================================================================
\art{
  \setsat{47}
  \begin{scope}[blend mode=screen]
     \fill[nightgrey, opacity=0.9, path fading=rglow] (8,4.5) circle (9);
     \fill[wred,  opacity=0.45, path fading=rglow] (4.4,4.6) circle (5);
     \fill[wcyan, opacity=0.40, path fading=rglow] (11.6,2.4) circle (3.5);\end{scope}
  \dust{80}
  \begin{scope}[blend mode=screen]
     \fill[wred, opacity=0.5, path fading=rglow] (4.4,4.6) ellipse (1.7 and 4.2);\end{scope}
  \shade[bottom color=wred!70!void, top color=wmag] (3.5,0.9) rectangle (5.3,8.3);
  \node[text=white, rotate=90, font=\footnotesize] at (4.4,4.6) {11{,}695{,}232{,}963 comparisons};
  \spark{(11.6,2.4)}{0.13}{wcyan}
  \node[text=wcyan, font=\footnotesize] at (11.6,3.4) {159{,}018 hash probes};
  \node[text=white] at (10.9,8.0) {\Large a list, walked};
  \node[text=white] at (10.9,7.0) {\Large vs.\ an index};
  \meter{47}
}

% =====================================================================
% 7. THREE ROTORS
% =====================================================================
\art{
  \setsat{57}
  \begin{scope}[blend mode=screen]
     \fill[nightgrey, opacity=0.9, path fading=rglow] (8,4.8) circle (9);\end{scope}
  \dust{70}
  \node[text=white] at (8,8.4) {\Large three rotors, one model --- the lesson travels};
  \draw[wamber, line width=1.5pt, opacity=0.5, -{Stealth[length=4mm]}]
     (4.4,6.0) .. controls (6,7.0) and (7,7.0) .. (7.6,6.2);
  \draw[wcyan, line width=1.5pt, opacity=0.6, -{Stealth[length=4mm]}]
     (8.6,6.2) .. controls (10,7.0) and (11,7.0) .. (11.8,6.0);
  \foreach \x/\c/\name/\sub in {%
      3.4/wamber/{Original C}/{the master},
      8.0/wviolet/{Hash-consed C}/{taught one trick},
      12.6/wcyan/{Rust}/{the lesson, fully}}{
     \fill[\c, opacity=0.5, path fading=rglow] (\x,5.4) circle (1.7);
     \shade[inner color=white, outer color=\c!55!void] (\x,5.4) circle (0.9);
     \draw[\c, line width=1.4pt] (\x,5.4) circle (0.9);
     \node[text=silver, font=\small\bfseries] at (\x,3.8) {\name};
     \node[text=\c, font=\footnotesize\itshape] at (\x,3.4) {\sub};}
  \node[text=silver, font=\footnotesize] at (3.4,2.7) {139\,s $\cdot$ 11.7\,B compares};
  \node[text=silver, font=\footnotesize] at (8.0,2.7) {1.5\,s $\cdot$ byte-identical};
  \node[text=silver, font=\footnotesize] at (12.6,2.7){0.1\,s $\cdot$ 1000$\times$};
  \node[text=silver, opacity=0.9] at (8,1.4)
     {\footnotesize the speed was never the language --- one small thing, done well};
  \meter{57}
}

% =====================================================================
% 8. THE BACK-PORT
% =====================================================================
\art{
  \setsat{66}
  \begin{scope}[blend mode=screen]
     \fill[nightgrey, opacity=0.9, path fading=rglow] (8,4.7) circle (8.5);
     \fill[wamber, opacity=0.28, path fading=rglow] (3.9,4.7) circle (3.2);
     \fill[wcyan,  opacity=0.28, path fading=rglow] (12.1,4.7) circle (3.2);\end{scope}
  \dust{70}
  \rays{(3.9,4.7)}{wgold}{15}\rays{(12.1,4.7)}{wcyan}{15}
  \shade[inner color=white, outer color=wamber!70!void] (3.9,4.7) circle (1.05);
  \node[text=void, font=\bfseries] at (3.9,4.7) {C};
  \shade[inner color=white, outer color=wcyan!70!void] (12.1,4.7) circle (1.05);
  \node[text=void, font=\bfseries] at (12.1,4.7) {C\,+\,hash};
  \draw[white, line width=6pt, opacity=0.08] (5.0,4.7) -- (11.0,4.7);
  \draw[wlime, line width=2.5pt, opacity=0.30,
        decorate, decoration={random steps, segment length=5pt, amplitude=2pt}] (5.0,4.7) -- (11.0,4.7);
  \draw[white, line width=1pt, opacity=0.85,
        decorate, decoration={random steps, segment length=5pt, amplitude=2pt}] (5.0,4.7) -- (11.0,4.7);
  \spark{(8,4.7)}{0.10}{wlime}
  \node[text=white] at (8,8.0) {\Large the same idea, carried back into C};
  \glowtext{(8,2.2)}{wlime}{\fontsize{38}{40}\selectfont\textbf{95$\times$} \normalsize\textcolor{silver}{, byte-identical}}
  \meter{66}
}

% =====================================================================
% 9. SYMBOLIC ARGV
% =====================================================================
\art{
  \setsat{76}
  \begin{scope}[blend mode=screen]
     \fill[nightgrey, opacity=0.9, path fading=rglow] (8,4.6) circle (9);
     \fill[weblue, opacity=0.4, path fading=rglow] (8,4.6) circle (7);\end{scope}
  \dust{70}
  \foreach \i/\open/\c in {0/0/silver,1/0/silver,2/1/wcyan,3/1/wmag,4/1/wlime,5/1/wamber}{
     \pgfmathsetmacro{\lx}{2.0+\i*2.35}
     \ifnum\open=1
        \fill[\c, opacity=0.45, path fading=rglow] (\lx,4.7) circle (1.5);
        \begin{scope}[blend mode=screen]
        \foreach \pk in {1,...,7}{
           \pgfmathsetmacro{\dy}{rnd*2.4}\pgfmathsetmacro{\dx}{(rnd-0.5)*0.7}
           \fill[\c, opacity={0.5-0.15*rnd}, path fading=rglow] (\lx+\dx,5.4+\dy) circle (0.07);}
        \end{scope}
     \fi
     \shade[top color=\c!75!void, bottom color=\c!30!void] (\lx-0.5,3.5) rectangle (\lx+0.5,4.7);
     \draw[\c, line width=1.5pt] (\lx-0.5,3.5) rectangle (\lx+0.5,4.7);
     \ifnum\open=1 \draw[\c, line width=1.6pt] (\lx-0.30,4.7) arc (200:20:0.36);
     \else \draw[\c, line width=1.6pt] (\lx-0.30,4.7) arc (180:0:0.30);\fi
     \fill[white] (\lx,4.1) circle (0.08);}
  \node[text=white] at (8,8.2) {\Large the command line: frozen $\rightarrow$ \textbf{open}};
  \node[text=silver] at (8,2.0) {\large the solver fills the open bytes};
  \meter{76}
}

% =====================================================================
% 10. THE X / Y SECRET
% =====================================================================
\art{
  \setsat{84}
  \begin{scope}[blend mode=screen]
     \fill[nightgrey, opacity=0.9, path fading=rglow] (8,5.0) circle (9);
     \fill[wmag, opacity=0.5, path fading=rglow] (8,5.0) circle (5.5);\end{scope}
  \dust{70}
  \foreach \dx/\ch/\c in {5.4/X/wcyan, 10.6/Y/wlime}{
     \fill[\c, opacity=0.45, path fading=rglow] (\dx,5.1) circle (2.1);
     \shade[inner color=\c!35!void, outer color=void] (\dx,5.1) circle (1.35);
     \draw[\c, line width=2pt] (\dx,5.1) circle (1.35);
     \foreach \t in {0,15,...,345}{\draw[\c,opacity=0.55,line width=1pt]
        ({\dx+1.16*cos(\t)},{5.1+1.16*sin(\t)}) -- ({\dx+1.35*cos(\t)},{5.1+1.35*sin(\t)});}
     \node[text=white] at (\dx,5.1) {\fontsize{40}{42}\selectfont\textbf{\ch}};}
  \rays{(8,5.1)}{wgold}{28}
  \begin{scope}[blend mode=screen]
  \foreach \t in {0,18,...,342}{
     \pgfmathsetmacro{\rr}{0.7+0.6*rnd}
     \draw[wamber,opacity=0.8,line width=1.2pt]
        ({8+0.3*cos(\t)},{5.1+0.3*sin(\t)}) -- ({8+\rr*cos(\t)},{5.1+\rr*sin(\t)});}
  \end{scope}
  \spark{(8,5.1)}{0.10}{wamber}
  \node[text=white] at (8,8.5) {\Large it crashes only on \textbf{X} and \textbf{Y} together};
  \node[text=wpink] at (8,1.2) {\large found by reason, not by typing};
  \meter{84}
}

% =====================================================================
% 11. 36 / 36
% =====================================================================
\art{
  \setsat{90}
  \begin{scope}[blend mode=screen]
     \fill[nightgrey, opacity=0.9, path fading=rglow] (8,5.4) circle (9);
     \fill[wgreen, opacity=0.40, path fading=rglow] (8,5.4) circle (7);\end{scope}
  \dust{60}\rays{(8,5.4)}{wgreen}{16}
  \foreach \r in {0,...,3}{\foreach \cc in {0,...,8}{
     \pgfmathsetmacro{\cx}{2.6+\cc*1.35}
     \pgfmathsetmacro{\cy}{7.0-\r*1.05}
     \fill[wgreen, opacity=0.45, path fading=rglow] (\cx,\cy) circle (0.42);
     \draw[white, line width=1.8pt, line cap=round]
        (\cx-0.18,\cy) -- (\cx-0.03,\cy-0.16) -- (\cx+0.24,\cy+0.22);}}
  \glowtext{(8,1.7)}{wgreen}{\fontsize{46}{48}\selectfont\textbf{36 / 36}}
  \node[text=silver, font=\footnotesize] at (8,0.8)
    {same bug, same step --- both tools, every benchmark};
  \meter{90}
}

% =====================================================================
% 12. THE VISUALIZER
% =====================================================================
\art{
  \setsat{95}
  \begin{scope}[blend mode=screen]
     \fill[nightgrey, opacity=0.9, path fading=rglow] (8,4.4) circle (9);
     \fill[wred,    opacity=0.40, path fading=rglow] (8,7.3) circle (3.5);
     \fill[wviolet, opacity=0.35, path fading=rglow] (8,4) circle (7);\end{scope}
  \dust{90}
  \coordinate (top) at (8,6.7);
  \coordinate (na) at (5.6,5.4); \coordinate (nb) at (10.4,5.4);
  \coordinate (nc) at (3.4,3.6); \coordinate (nd) at (6.8,3.6);
  \coordinate (ne) at (9.2,3.6); \coordinate (nf) at (12.2,3.6);
  \coordinate (ng) at (4.8,1.9); \coordinate (nh) at (7.6,1.9);
  \foreach \a/\b/\c in {top/na/wred, top/nb/wred, na/nc/wcyan, na/nd/wviolet,
        nb/ne/wcyan, nb/nf/wlime, nd/ng/wamber, nd/nh/wviolet}{\neon{(\a)}{(\b)}{\c}}
  \rays{(8,7.4)}{wred}{20}
  \node[star, star points=8, star point ratio=0.55, draw=wred, fill=wred!70!void,
        line width=1.3pt, minimum size=1.7cm] at (8,7.4) {};
  \node[text=white, font=\tiny\bfseries] at (8,7.4) {VIOLATED};
  \foreach \n/\c in {na/wcyan, nb/wteal, nc/weblue, nd/wviolet, ne/weblue,
        nf/wlime, ng/wamber, nh/wpink}{\spark{(\n)}{0.085}{\c}}
  \node[text=white] at (8,0.8) {\Large the unreadable answer, made watchable};
  \meter{95}
}

% =====================================================================
% 13. THE LESSON
% =====================================================================
\art{
  \setsat{98}
  \begin{scope}[blend mode=screen]
     \fill[nightgrey, opacity=0.9, path fading=rglow] (8,5) circle (9);
     \fill[wviolet, opacity=0.4, path fading=rglow] (8,5) circle (7);
     \fill[wcyan, opacity=0.28, path fading=rglow] (8,5.2) circle (3.2);\end{scope}
  \dust{110}\rays{(8,5.2)}{wcyan}{16}
  \shade[inner color=white, outer color=weblue!55!void]
     (8,7.8) -- (5.8,6.6) -- (5.8,3.9)
     .. controls (5.8,2.6) and (7.0,2.1) .. (8,1.6)
     .. controls (9.0,2.1) and (10.2,2.6) .. (10.2,3.9)
     -- (10.2,6.6) -- cycle;
  \draw[wcyan, line width=2pt]
     (8,7.8) -- (5.8,6.6) -- (5.8,3.9)
     .. controls (5.8,2.6) and (7.0,2.1) .. (8,1.6)
     .. controls (9.0,2.1) and (10.2,2.6) .. (10.2,3.9)
     -- (10.2,6.6) -- cycle;
  \begin{scope}[blend mode=screen]
     \fill[white, opacity=0.8, path fading=rglow] (7.7,5.2) circle (1.3);\end{scope}
  \draw[void, line width=3pt, line cap=round, line join=round]
     (7.2,4.8) -- (7.75,3.9) -- (8.9,5.9);
  \node[text=white] at (8,0.85) {\Large the hard part was proving it still tells the truth};
  \meter{98}
}

% =====================================================================
% 14. THANK YOU  -- full colour
% =====================================================================
\art{
  \setsat{100}
  \nebula
  \begin{scope}[blend mode=screen]
  \foreach \x/\c/\o in {2.5/wgreen/0.22, 4/wteal/0.18, 5.8/wcyan/0.20,
        9.5/wviolet/0.18, 11.5/wmag/0.16, 13.5/weblue/0.18}{
     \fill[\c, opacity=\o, path fading=fadeup] (\x-0.7,4.5) rectangle (\x+0.7,9);}
  \end{scope}
  \dust{240}\bokeh{55}
  \coordinate (r1) at (4.6,6.3);\coordinate (r2) at (8.2,7.0);\coordinate (r3) at (11.4,5.6);
  \neon{(r1)}{(r2)}{wcyan}\neon{(r2)}{(r3)}{wmag}
  \spark{(r1)}{0.09}{wcyan}\spark{(r2)}{0.09}{wpink}\spark{(r3)}{0.09}{wlime}
  \begin{scope}[blend mode=screen]
     \fill[wcyan, opacity=0.40, path fading=rglow] (8,3.4) ellipse (5.5 and 1.5);\end{scope}
  \node[text=white] at (8,3.5) {\fontsize{48}{52}\selectfont\textbf{Thank you}};
  \node[text=silver, opacity=0.75] at (8,2.3)
    {\footnotesize github.com/jaspek/rotor-rust \;$\cdot$\; rotor-c-hashcons
     \;$\cdot$\; jaspek.github.io/rotor-rust};
  \meter{100}
}

\end{document}
